% !TEX program = xelatex
\documentclass[a4paper]{article}
\usepackage[UTF8]{ctex}
\setCJKmainfont{SimSun}[BoldFont=SimHei,ItalicFont=KaiTi]
\setCJKsansfont{SimHei}
\usepackage[top=4mm,bottom=4mm,left=4mm,right=4mm]{geometry}
\usepackage{amsmath,amssymb,bm}
\usepackage{multicol}
\usepackage{graphicx}
\usepackage[dvipsnames]{xcolor}
\usepackage{enumitem}
\usepackage{array,tabularx}
\pagestyle{empty}
\setlength{\parindent}{0pt}
\setlength{\parskip}{0.2pt}
\setlength{\columnsep}{2.8mm}
\setlength{\columnseprule}{0.15pt}
\renewcommand{\baselinestretch}{0.80}
\setlist{nosep,leftmargin=7pt,labelsep=1.5pt,itemsep=0pt,parsep=0pt,topsep=0pt}
\definecolor{h1}{HTML}{0D47A1}
\definecolor{h2}{HTML}{1565C0}
\definecolor{h3}{HTML}{E65100}
\definecolor{warnc}{HTML}{B71C1C}
\definecolor{tip}{HTML}{1B5E20}
\definecolor{bg}{HTML}{EEF2FF}
\definecolor{bgw}{HTML}{FFF3E0}
\newcommand{\chap}[1]{\vspace{1pt}{\fontsize{7.0}{7.3}\selectfont\bfseries\color{white}\colorbox{h1}{\strut\;#1\;}}\par\vspace{0.4pt}}
\newcommand{\sect}[1]{\vspace{0.5pt}{\fontsize{5.9}{6.2}\selectfont\bfseries\color{h2}$\blacktriangleright$ #1}\par}
\newcommand{\exm}[1]{{\fontsize{5.45}{5.8}\selectfont\bfseries\color{tip}$\diamond$ #1}\par}
\newcommand{\key}[1]{{\bfseries\color{h3}#1}}
\newcommand{\warn}[1]{{\bfseries\color{warnc}#1}}
\newcommand{\fml}[1]{\colorbox{bg}{$\displaystyle #1$}}
\newcommand{\wfml}[1]{\colorbox{bgw}{$\displaystyle #1$}}
\newcommand{\miniimg}[2]{\begin{center}\includegraphics[width=#1\columnwidth]{#2}\end{center}\vspace{-3pt}}
\begin{document}
\fontsize{5.55}{6.15}\selectfont

% ==================== PAGE 1 ====================
\begin{multicols}{3}
\chap{P1 总判题地图：先选模型}
\miniimg{0.95}{assets/img08_models.jpg}
\sect{考试资料策略}
可带：指定教材+计算器+3张A4双面。教材查长定义，本表放\key{题型识别、公式、答案骨架、易错点}。目标不是背诵，而是看到题能定位模型。

\sect{题干关键词定位}
\renewcommand{\arraystretch}{0.76}\begin{tabular}{|p{0.30\linewidth}|p{0.61\linewidth}|}
\hline 静止/水深/闸门/U管 & 静水压：$p=p_0+\rho gh$；平面/曲面总压力\\ \hline
喷管/Ma/背压/激波 & P3/P4：等熵、临界、激波、膨胀波\\ \hline
势函数/流函数/圆柱 & P2：$\varphi,\psi$叠加+伯努利\\ \hline
黏度/Re/管流/泵 & P5：N-S简化、H-P、管路损失\\ \hline
平板/边界层/阻力 & P6：$C_f,\delta,D$\\ \hline
弯管/喷流/受力 & 控制体动量：$\sum F=\dot m(V_2-V_1)$\\ \hline
\end{tabular}\renewcommand{\arraystretch}{1}

\sect{常数与基本量}
$g=9.81$；水$\rho=1000$；空气$\gamma=1.4,R=287$；$1\rm atm=101325Pa=10.33mH_2O$。\\
\fml{\nu=\mu/\rho,\quad Re=VL/\nu=\rho VL/\mu,\quad Ma=V/a,\quad a=\sqrt{\gamma RT}}

\sect{连续性}
不可压：$\nabla\cdot\bm V=0$；一维定常：$A_1V_1=A_2V_2=Q$。非均匀速度必须积分：$Q=\int_A v\,dA$。圆管抛物线层流：$\bar v=v_{\max}/2$。

\sect{控制体动量模板}
\fml{\sum F_x=\dot m(V_{2x}-V_{1x}),\quad \sum F_y=\dot m(V_{2y}-V_{1y})}, $\dot m=\rho Q$。\\
步骤：画CV；标速度方向；列压力/重力/支反力；动量写\key{流出-流入}；物体受力与流体受力反向。\\
\exm{弯管}
$x:p_1A_1-p_2A_2\cos\alpha-F_{Rx}=\rho Q(V_2\cos\alpha-V_1)$；\\
$y:-p_2A_2\sin\alpha+F_{Ry}=\rho QV_2\sin\alpha$。
\exm{射流冲板}
垂直板：$F=\rho QV$。斜板/分流：分$x,y$列动量，未给损失时各支速度近似$V$。

\sect{伯努利/机械能}
理想同流线：\fml{p/\rho g+V^2/(2g)+z=C}。\\
含损失/泵/水轮机：\\
\fml{\frac{p_1}{\rho g}+\frac{\alpha_1V_1^2}{2g}+z_1+H_p=\frac{p_2}{\rho g}+\frac{\alpha_2V_2^2}{2g}+z_2+h_f+\sum h_\zeta+H_T}
\warn{泵加左，水轮机输出加右；开口液面表压为0。}

\sect{静水压旧题急救}
平面：$F=\rho gh_CA$，$y_D=y_C+J_C/(y_CA)$。\\
曲面：$F_x=\rho gh_{Cx}A_x$；$F_y=\rho gV_{\text{压}}$；圆柱面合力过圆心。\\
U管爬楼梯：向下$+\rho gh$，向上$-\rho gh$，同一静止连通液体同高等压。

\sect{概念判断}
亚声速椭圆型、全局影响；超声速双曲型、马赫锥内传播。激波绝热非等熵，$T_0$不变、$p_0$下降。膨胀波近似等熵。逆压梯度易分离。
\end{multicols}
\newpage

% ==================== PAGE 2 ====================
\begin{multicols}{3}
\chap{P2 势流叠加：理想不可压无旋}
\miniimg{0.95}{assets/img01_potential.jpg}
\sect{基本条件}
无旋：$\nabla\times\bm V=0\Rightarrow \bm V=\nabla\varphi$。二维不可压可设流函数$\psi$。\\
\fml{\nabla^2\varphi=0,\quad \nabla^2\psi=0}
直角：$u=\varphi_x=\psi_y,\;v=\varphi_y=-\psi_x$。\\
极坐标：\fml{v_r=\varphi_r=\frac1r\psi_\theta,\quad v_\theta=\frac1r\varphi_\theta=-\psi_r}

\sect{基本流表}
\renewcommand{\arraystretch}{0.76}\begin{tabular}{|p{0.18\linewidth}|p{0.35\linewidth}|p{0.35\linewidth}|}
\hline 流动 & $\varphi$ & $\psi$\\ \hline
均匀流 & $Ur\cos\theta$ & $Ur\sin\theta$\\ \hline
源 & $\frac{Q}{2\pi}\ln r$ & $\frac{Q}{2\pi}\theta$\\ \hline
汇 & $-\frac{Q}{2\pi}\ln r$ & $-\frac{Q}{2\pi}\theta$\\ \hline
点涡 & $\frac{\Gamma}{2\pi}\theta$ & $-\frac{\Gamma}{2\pi}\ln r$\\ \hline
偶极子 & $\frac{M\cos\theta}{r}$ & $-\frac{M\sin\theta}{r}$\\ \hline
\end{tabular}\renewcommand{\arraystretch}{1}
\warn{叠加法：总$\varphi,\psi$直接相加，速度也相加。}

\sect{给速度求$\varphi,\psi$}
有$\varphi$先验：$v_x-u_y=0$；有$\psi$先验：$u_x+v_y=0$。\\
$\varphi_x=u\Rightarrow \varphi=\int u dx+f(y)$，再用$\varphi_y=v$定$f$。\\
$\psi_y=u\Rightarrow \psi=\int u dy+g(x)$，再用$-\psi_x=v$定$g$。

\sect{圆柱绕流}
\fml{\varphi=U(r+a^2/r)\cos\theta,\quad \psi=U(r-a^2/r)\sin\theta}\\
$v_r=U(1-a^2/r^2)\cos\theta$，$v_\theta=-U(1+a^2/r^2)\sin\theta$。\\
壁面$r=a$：$v_r=0,\;v_\theta=-2U\sin\theta$；驻点$\theta=0,\pi$；最大速$2U$在$90^\circ,270^\circ$。\\
压强：\fml{p=p_\infty+\frac12\rho(U^2-V^2)}。壁面$p=p_\infty+\frac12\rho U^2(1-4\sin^2\theta)$。

\sect{有环量/升力}
圆柱+点涡：$v_\theta=-2U\sin\theta+\Gamma/(2\pi a)$。升力：\fml{L'=\rho U\Gamma}。理想无粘绕流阻力$D=0$。

\sect{镜像法/边角}
固壁为流线，要求法向速度为0。源/汇靠直壁：同号镜像；点涡靠直壁：反号镜像。半平面源汇常见$2\pi\to\pi$。\\
边角：若$\varphi=Ar^n\cos n\theta,\psi=Ar^n\sin n\theta$，壁面$\psi=0$，夹角$\alpha$满足$n=\pi/\alpha$。$V\sim r^{n-1}$。

\sect{答案句}
“该流动不可压无旋，故$\varphi,\psi$满足Laplace方程。由叠加法写总流函数，速度由极坐标关系求得，固壁为$\psi=C$。压强由伯努利求。”
\end{multicols}
\newpage

% ==================== PAGE 3 ====================
\begin{multicols}{3}
\chap{P3 等熵流 + Laval 喷管}
\miniimg{0.95}{assets/img04_laval.jpg}
\sect{基本公式}
\fml{p=\rho RT,\quad a=\sqrt{\gamma RT},\quad Ma=V/a}
\fml{\frac{T_0}{T}=1+\frac{\gamma-1}{2}M^2}
\fml{\frac{p_0}{p}=\left(1+\frac{\gamma-1}{2}M^2\right)^{\gamma/(\gamma-1)}}
\fml{\frac{\rho_0}{\rho}=\left(1+\frac{\gamma-1}{2}M^2\right)^{1/(\gamma-1)}}\\
空气：$T_0/T=1+0.2M^2$，$p_0/p=(1+0.2M^2)^{3.5}$，$\rho_0/\rho=(1+0.2M^2)^{2.5}$。

\sect{临界/阻塞}
$M=1,A=A^*$。\\
\fml{T^*/T_0=0.833,\quad p^*/p_0=0.528,\quad \rho^*/\rho_0=0.634}\\
\warn{0.528只表示临界静压比，不是所有出口/背压比。}

\sect{面积--马赫数}
\fml{\frac{A}{A^*}=\frac1M\left[\frac{2}{\gamma+1}\left(1+\frac{\gamma-1}{2}M^2\right)\right]^{\frac{\gamma+1}{2(\gamma-1)}}}\\
$A/A^*>1$通常有亚声速/超声速两支。收缩管最多到$M=1$；Laval扩张段可超声速。

\sect{质量流量}
\fml{\dot m=\rho VA=ApM\sqrt{\gamma/(RT)}}\\
阻塞最大流量在$M=1$，空气：\fml{\dot m_{\max}=0.0404A^*p_0/\sqrt{T_0}}。

\sect{Laval 背压状态}
\renewcommand{\arraystretch}{0.76}\begin{tabular}{|p{0.30\linewidth}|p{0.61\linewidth}|}
\hline 背压高 & 全管亚声速，$p_e=p_b$，未阻塞\\ \hline
喉部$M=1$ & 阻塞，$\dot m$最大\\ \hline
扩张段正激波 & 波前超声速，波后亚声速，$p_0$降\\ \hline
设计状态 & 全管等熵，$p_e=p_b$\\ \hline
$p_e<p_b$ & 过膨胀，出口外压缩/斜激波\\ \hline
$p_e>p_b$ & 欠膨胀，出口外膨胀波\\ \hline
\end{tabular}\renewcommand{\arraystretch}{1}

\sect{收缩喷管题}
给$p_0,T_0,p_b,A_e$。若$p_b/p_0>0.528$：未阻塞，$p_e=p_b$，用$p_0/p_e$求$M_e$，再求$T_e,\rho_e,V_e,\dot m=\rho_eV_eA_e$。\\
若$p_b/p_0\le0.528$：阻塞，出口/喉部$M=1$，$\dot m=0.0404A_ep_0/\sqrt{T_0}$。

\sect{Laval面积/流量题}
给$A^*$与出口$p_e$或$T_e$，无激波。用等熵求$M_e$；若取超声速支，则喉部$M=1,A_t=A^*$。出口面积由$A_e/A^*$求；流量可用喉部最大流量或出口$\rho VA$。

\sect{文丘里质量流量计}
给最大流量和最低喉压$p_t$。先由入口$M_1$求$p_0,T_0$，再用$p_0/p_t=(1+0.2M_t^2)^{3.5}$求$M_t$。若$M_t<1$，不可用0.528，继续算$T_t,\rho_t,V_t,A_t$。
\end{multicols}
\newpage

% ==================== PAGE 4 ====================
\begin{multicols}{3}
\chap{P4 激波 + 膨胀波}
\miniimg{0.82}{assets/img02_shock.jpg}
\miniimg{0.82}{assets/img03_pm.jpg}
\sect{总判断}
超声速遇压缩角/凹角/迎风面$\Rightarrow$激波；遇扩张角/凸角$\Rightarrow$膨胀波；偏转角太大$\Rightarrow$脱体激波。激波：$p,T,\rho\uparrow$，$V,M\downarrow$，$T_0$不变，$p_0$下降。

\sect{正激波}
\fml{M_2^2=\frac{1+\frac{\gamma-1}{2}M_1^2}{\gamma M_1^2-\frac{\gamma-1}{2}}}
空气：$M_2^2=(1+0.2M_1^2)/(1.4M_1^2-0.2)$。\\
\fml{\frac{p_2}{p_1}=1+\frac{2\gamma}{\gamma+1}(M_1^2-1)}
\fml{\frac{\rho_2}{\rho_1}=\frac{(\gamma+1)M_1^2}{(\gamma-1)M_1^2+2}}
\fml{\frac{T_2}{T_1}=(p_2/p_1)/(\rho_2/\rho_1)}\\
$p_{02}/p_{01}=(p_2/p_1)(p_{02}/p_2)/(p_{01}/p_1)$，用等熵式算$p_0/p$。

\sect{斜激波流程}
$M_{n1}=M_1\sin\beta$。本质：法向分量过正激波，切向分量不变。\\
已知$M_1,\theta$：查$\theta-\beta-M$表得$\beta$；算$M_{n1}$；用正激波公式得$p_2/p_1,T_2/T_1,M_{n2}$；$M_2=M_{n2}/\sin(\beta-\theta)$。工程常取弱激波解。

\sect{膨胀波}
\fml{\nu(M)=\sqrt{\frac{\gamma+1}{\gamma-1}}\tan^{-1}\sqrt{\frac{\gamma-1}{\gamma+1}(M^2-1)}-\tan^{-1}\sqrt{M^2-1}}\\
空气：$\nu(M)=\sqrt6\tan^{-1}\sqrt{(M^2-1)/6}-\tan^{-1}\sqrt{M^2-1}$。\\
膨胀角：\fml{\delta=\nu(M_2)-\nu(M_1)}。给$\delta$：$\nu(M_2)=\nu(M_1)+\delta$。

\sect{7.19 管外膨胀模板}
给$A_e/A_t,p_b,p_0$。1 由$A_e/A^*$取超声速支求$M_e$。2 等熵求$p_e=p_0/(1+0.2M_e^2)^{3.5}$。3 若$p_e>p_b$，出口外继续膨胀。4 由$p_0/p_b$求最终$M_b$。5 $\alpha=\nu(M_b)-\nu(M_e)$。

\sect{7.20 管内正激波模板}
给激波位置$A_s/A_t$。1 激波前由$A_s/A^*$超声速支求$M_1$。2 正激波求$M_2,p_2/p_1,p_{02}/p_{01}$。3 激波后以新总压$p_{02}$亚声速等熵到出口。

\sect{易错}
知道$M_2$不能随便算全部：激波后静量接激波前$p_1,T_1,\rho_1$。$p_{01}$是波前滞止压，$p_{02}$是波后滞止压，$p_{02}<p_{01}$。膨胀波不用正激波公式。
\end{multicols}
\newpage

% ==================== PAGE 5 ====================
\begin{multicols}{3}
\chap{P5 粘性管流：N-S、H-P、Couette、管路}
\miniimg{0.92}{assets/img05_viscous.jpg}
\sect{N-S 简化}
不可压牛顿流体：$\rho D\bm V/Dt=-\nabla p+\mu\nabla^2\bm V+\rho\bm g$。定常$\partial_t=0$；单向流只保留$u$；充分发展$\partial u/\partial x=0$；平行流惯性项常为0。\\
\warn{大题先写假设，再积分两次，再代边界条件。}

\sect{H-P 圆管层流}
\fml{u(r)=\frac{1}{4\mu}\left(-\frac{dp}{dx}\right)(R^2-r^2)}
\fml{u_{\max}=2\bar u,\quad Q=\frac{\pi R^4}{8\mu}\left(-\frac{dp}{dx}\right)}
\fml{\Delta p=\frac{8\mu LQ}{\pi R^4},\quad \lambda=64/Re}
壁剪：$\tau_w=R\Delta p/(2L)=8\mu\bar u/d$。

\sect{Couette 流}
无压差：$u=Uy/h,\;\tau=\mu U/h$。\\
有压力梯度/重力分量：\fml{\mu d^2u/dy^2=dp/dx-\rho g_x}，积分两次。\\
倾斜平板若$u(0)=0,u(H)=-U_0,u(H/2)=0$，且$\mu u''+\rho g\sin\theta=0$，则$U_0=\rho g\sin\theta H^2/(4\mu)$。

\sect{双层 Couette}
\fml{\tau=\frac{U}{h_1/\mu_1+h_2/\mu_2}=\frac{U\mu_1\mu_2}{\mu_2h_1+\mu_1h_2}}
$u_1=\tau y/\mu_1$；$u_2=u_1(h_1)+\tau(y-h_1)/\mu_2$。

\sect{管路损失}
\fml{h_f=\lambda\frac{L}{d}\frac{V^2}{2g},\quad h_\zeta=\zeta\frac{V^2}{2g}}
粗糙度只通过$\varepsilon/d$影响$\lambda$；层流$\lambda=64/Re$与粗糙度无关。\\
吸水泵：$h=(p_a-p_B)/\rho g-V^2/2g-(\lambda L/d+\sum\zeta)V^2/2g$。\\
水轮机：$H_T=z_1-z_2+(p_1-p_2)/\rho g+(V_1^2-V_2^2)/2g-h_f-\sum h_\zeta$，$P=\rho gQH_T$。

\sect{常见题模板}
\exm{泵最大安装高度}
算$Q,V,Re,\varepsilon/d\to\lambda$；入口压力要求$p_B\ge p_v+\Delta p_{\text{safe}}$；代入求$h_{\max}$。
\exm{水轮机功率}
各段算$Re,\lambda,h_f$；能量方程含$H_T$；$P=\rho gQH_T$。
\exm{滑动轴承/油膜}
$du/dy\approx V/t$；$\tau=\mu V/t$；$F=\tau A$；$P=FV$。
\end{multicols}
\newpage

% ==================== PAGE 6 ====================
\begin{multicols}{3}
\chap{P6 边界层、阻力、最后急救}
\miniimg{0.94}{assets/img06_bl.jpg}
\sect{平板边界层}
$Re_x=Ux/\nu,\;Re_L=UL/\nu$。层流：\\
\fml{\delta=5x/\sqrt{Re_x},\quad \bar C_f=1.328/\sqrt{Re_L}}\\
湍流：\fml{\bar C_f=0.074/Re_L^{1/5}}；转捩：\fml{\bar C_f=0.074/Re_L^{1/5}-1742/Re_L}。\\
\fml{D=\bar C_f\frac12\rho U^2A,\quad P=DU}

\sect{位移/动量厚度}
\fml{\delta^*=\int_0^\delta(1-u/U)dy,\quad \theta=\int_0^\delta(u/U)(1-u/U)dy}
\fml{\tau_w/(\rho U^2)=d\theta/dx}
给$u/U=f(\eta),\eta=y/\delta$：把$dy=\delta d\eta$代入积分。

\sect{书10.2 平板油流}
给$\mu,\rho,U,L,b$。先$\nu=\mu/\rho$，$Re_L=UL/\nu$。若层流：$\delta_{\max}=5L/\sqrt{Re_L}$，两面阻力$D=\bar C_f(0.5\rho U^2)(2Lb)$。\\
\warn{这是平板边界层，不是管流壁律。}

\sect{湍流近壁三层}
$u_\tau=\sqrt{\tau_w/\rho}$，$y^+=yu_\tau/\nu$，$u^+=\bar u/u_\tau$。线性底层$y^+<5,u^+=y^+$；缓冲层$5<y^+<30$；湍流核心$u^+=2.5\ln y^+ +5.5$。圆管平均：$\bar u/u_\tau=2.5\ln(Ru_\tau/\nu)+1.75$。

\sect{分离与阻力}
逆压梯度$dp/dx>0$使边界层减速增厚；$\partial u/\partial y|_w=0$为分离临界。摩擦阻力来自$\tau_w$；压差阻力来自分离尾流。

\sect{最后急救题库}
\renewcommand{\arraystretch}{0.76}\begin{tabular}{|p{0.30\linewidth}|p{0.61\linewidth}|}
\hline 速度场连续 & 算$u_x+v_y+w_z$\\ \hline
有旋 & $\Omega_z=v_x-u_y$，$\omega_z=\Omega_z/2$\\ \hline
流线 & $dx/u=dy/v$，非定常先代$t=t_0$\\ \hline
势/流函数 & 先验无旋/不可压，再积分\\ \hline
小孔出流 & $V=\sqrt{2gH}$，$Q=\mu A\sqrt{2gH}$\\ \hline
皮托管 & $V=\sqrt{2(p_0-p)/\rho}$\\ \hline
U管差压 & $\Delta p=(\rho_m-\rho)g\Delta h$\\ \hline
维数分析 & $n$个量、$k$个基本量纲，得$n-k$个$\pi$数\\ \hline
\end{tabular}\renewcommand{\arraystretch}{1}

\sect{不会也要写的步骤分}
先写“该题属于……模型”，再写控制方程、边界条件、代入顺序。比如：“定常一维绝热等熵流，由$M$求滞止量，由面积比求另一截面$M$，再由等熵式求$p,T,\rho,V$。”这能拿框架分。
\end{multicols}
\end{document}
